% ============================================================
%  Taller 2 — Segmentación por Umbralización en Imágenes Médicas
%  Pontificia Universidad Javeriana
%  Procesamiento de Imágenes Médicas
%
%  Compilar con: pdflatex taller2_segmentacion.tex
%  (dos pasadas para referencias cruzadas)
%
%  Estructura de carpetas esperada desde la raíz del proyecto:
%    output/comparisons/   ← imágenes *_comparison.png
% ============================================================

\documentclass[12pt, a4paper]{article}

% ── Codificación y lenguaje ────────────────────────────────
\usepackage[utf8]{inputenc}
\usepackage[T1]{fontenc}
\usepackage[spanish, es-tabla]{babel}

% ── Tipografía y geometría ─────────────────────────────────
\usepackage{lmodern}
\usepackage[margin=2.5cm]{geometry}
\usepackage{microtype}
\usepackage{parskip}
\setlength{\parindent}{0pt}
\setlength{\parskip}{6pt}

% ── Matemáticas ────────────────────────────────────────────
\usepackage{amsmath, amssymb, amsthm}

% ── Figuras y subfiguras ───────────────────────────────────
\usepackage{graphicx}
\usepackage{subcaption}
\usepackage{float}
\usepackage[export]{adjustbox}

% ── Tablas ────────────────────────────────────────────────
\usepackage{booktabs}
\usepackage{tabularx}
\usepackage{multirow}
\usepackage{xcolor}
\usepackage{colortbl}
\usepackage{array}

% ── Hipervínculos ─────────────────────────────────────────
\usepackage[colorlinks=true, linkcolor=blue!70!black,
            citecolor=green!50!black, urlcolor=blue!70!black]{hyperref}

% ── Encabezados y pies de página ──────────────────────────
\usepackage{fancyhdr}
\pagestyle{fancy}
\fancyhf{}
\fancyhead[L]{\small Pontificia Universidad Javeriana $\mid$ Procesamiento de Imágenes Médicas}
\fancyhead[R]{\small Taller 2 — Segmentación por Umbralización}
\fancyfoot[C]{\thepage}
\renewcommand{\headrulewidth}{0.4pt}

% ── Cajas de información ───────────────────────────────────
\usepackage{tcolorbox}
\tcbuselibrary{skins, breakable}
\tcbset{
  infostyle/.style={
    colback=blue!5, colframe=blue!40!black,
    fonttitle=\bfseries, breakable, enhanced
  },
  warnstyle/.style={
    colback=orange!8, colframe=orange!70!black,
    fonttitle=\bfseries, breakable, enhanced
  },
  codestyle/.style={
    colback=gray!6, colframe=gray!50,
    fonttitle=\bfseries, breakable, enhanced
  }
}

% ── Títulos de sección personalizados ─────────────────────
\usepackage{titlesec}
\titleformat{\section}{\large\bfseries\color{blue!70!black}}{}{0em}
  {\thesection\quad}[\titlerule]
\titleformat{\subsection}{\normalsize\bfseries\color{blue!50!black}}
  {\thesubsection\quad}{0em}{}
\titleformat{\subsubsection}{\normalsize\itshape\bfseries}
  {\thesubsubsection\quad}{0em}{}

% ── Rutas base de imágenes ────────────────────────────────
% El .tex está en: umbrales-ITK/
% Las imágenes están en: umbrales-ITK/output/
\graphicspath{
  {output/}
}

% ── Metadatos del documento ───────────────────────────────
\title{
  \vspace{-1cm}
  {\large Pontificia Universidad Javeriana}\\[0.3em]
  {\large Procesamiento de Imágenes Médicas}\\[1em]
  \rule{\linewidth}{1.5pt}\\[0.5em]
  {\LARGE\bfseries Taller 2 — Segmentación por Umbralización\\[0.2em]
   en Imágenes Médicas Multimodales}\\[0.3em]
  \rule{\linewidth}{1.5pt}
}
\author{
  Abel Albuez\quad$\mid$\quad Victoria Acero\quad$\mid$\quad Santiago Gil\\[0.4em]
  \small\texttt{aa-albuezs@javeriana.edu.co}
  $\;\cdot\;$
  \texttt{jv.acero@javeriana.edu.co}
  $\;\cdot\;$
  \texttt{santiago\_gil@javeriana.edu.co}
}
\date{Bogotá D.C., Colombia — Marzo 2026}

% ============================================================
\begin{document}
% ============================================================

\maketitle
\thispagestyle{fancy}

% ── Resumen ───────────────────────────────────────────────
\begin{tcolorbox}[infostyle, title={Resumen}]
Este documento reporta la implementación y análisis experimental de tres métodos clásicos
de umbralización automática aplicados a imágenes médicas de tres modalidades distintas:
tomografía computarizada de tórax (CT), resonancia magnética cerebral con contraste (MRI
T1c) y ultrasonido endorrectal de próstata (US). Los métodos evaluados son el umbral de
\textbf{Otsu}, el umbral de \textbf{Huang} y el umbral de \textbf{Triangle}, implementados
con la biblioteca \emph{Insight Toolkit} (ITK) en Python. El análisis examina el
comportamiento de cada método ante histogramas de diferente distribución y discute sus
implicaciones clínicas en el contexto de la segmentación de imágenes médicas.
\end{tcolorbox}

\vspace{0.5em}
\tableofcontents
\newpage

% ============================================================
\section{Introducción}
% ============================================================

La segmentación de imágenes médicas consiste en particionar una imagen en regiones
homogéneas respecto a una o más características, con el objetivo de aislar estructuras
anatómicas o patológicas de interés clínico \cite{bankman2009}. Entre los métodos más
simples y ampliamente utilizados en la práctica clínica se encuentra la
\textbf{umbralización}, que transforma la imagen original en una máscara binaria a partir
de un valor de intensidad umbral $T$: los vóxeles con intensidad inferior a $T$ se asignan
a una clase (fondo o tejido de interés) y los superiores a la otra.

La eficacia de la umbralización depende fundamentalmente de la estructura del histograma
de la imagen. El caso ideal asume un histograma \textbf{bimodal}: dos picos diferenciados
separados por un valle, donde cada pico representa una clase y el valle indica la ubicación
natural del umbral \cite{bankman2009}. Sin embargo, las imágenes médicas presentan
frecuentemente histogramas complejos, con múltiples modos, picos asimétricos o intensidades
superpuestas entre clases, lo que hace necesario el uso de métodos de umbralización
automática con criterios estadísticos robustos.

El objetivo de este taller es implementar y comparar tres métodos de umbralización
automática sobre tres modalidades de imagen médica, evaluando cualitativamente la calidad
de la segmentación resultante y discutiendo la adecuación de cada método según las
características del histograma de cada modalidad.

\begin{table}[H]
\centering
\renewcommand{\arraystretch}{1.3}
\begin{tabularx}{\textwidth}{llX}
\toprule
\textbf{Método} & \textbf{Filtro ITK} & \textbf{Criterio de optimización} \\
\midrule
Otsu   & \texttt{itk::OtsuThresholdImageFilter}     & Maximizar la varianza entre clases \\
Huang  & \texttt{itk::HuangThresholdImageFilter}    & Minimizar la entropía difusa (\textit{fuzzy}) \\
Triangle & \texttt{itk::TriangleThresholdImageFilter} & Maximizar la distancia geométrica al histograma \\
\bottomrule
\end{tabularx}
\caption{Métodos de umbralización automática implementados en este taller.}
\end{table}

% ============================================================
\section{Marco Teórico}
% ============================================================

\subsection{Umbralización global}

La umbralización global es el método de segmentación más simple: asume que la imagen
posee un histograma bimodal y que un único umbral $T$ permite separar los dos modos.
La imagen segmentada $g(x,y)$ se obtiene como:

\begin{equation}
  g(x,y) =
  \begin{cases}
    1 & \text{si } f(x,y) \geq T \\
    0 & \text{si } f(x,y) < T
  \end{cases}
  \label{eq:threshold}
\end{equation}

donde $f(x,y)$ es la intensidad del vóxel en la imagen original y los valores 0 y 1
representan el fondo y el objeto de interés, respectivamente. La elección de $T$ determina
completamente la calidad de la segmentación; métodos automáticos de estimación de $T$ son
preferibles a la selección manual para garantizar reproducibilidad.

\subsection{Método de Otsu}

El método de Otsu \cite{otsu1979} estima el umbral óptimo $T^*$ como aquel que
\textbf{maximiza la varianza entre clases} $\sigma_B^2(T)$, equivalente a minimizar la
varianza intra-clase ponderada. Para un umbral $T$ que divide el histograma normalizado
$p_i$ en dos clases $C_0 = \{0, \ldots, T-1\}$ y $C_1 = \{T, \ldots, L-1\}$:

\begin{equation}
  \sigma_B^2(T) = \omega_0(T)\,\omega_1(T)\,\bigl[\mu_0(T) - \mu_1(T)\bigr]^2
  \label{eq:otsu}
\end{equation}

donde $\omega_k$ es la probabilidad acumulada de la clase $k$ y $\mu_k$ su media de
intensidades. El umbral óptimo es:

\begin{equation}
  T^* = \arg\max_{T}\; \sigma_B^2(T)
\end{equation}

Otsu produce resultados óptimos cuando las dos clases tienen distribuciones gaussianas
bien separadas y frecuencias de aparición similares. Su rendimiento se degrada cuando las
clases tienen tamaños muy diferentes (una clase domina el histograma) o cuando los modos
se solapan considerablemente.

\subsection{Método de Huang}

El método de Huang y Wang \cite{huang1995} plantea el problema de umbralización desde
la perspectiva de la \textbf{lógica difusa} (\textit{fuzzy logic}). En lugar de una
asignación binaria rígida, cada vóxel tiene un grado de pertenencia $\mu(x)$ a cada clase,
definido por una función de membresía sigmoidal. El umbral óptimo es aquel que minimiza
la \textbf{entropía difusa total} $E(T)$ de la imagen:

\begin{equation}
  E(T) = -\sum_{i=0}^{L-1} p_i \left[ \mu_i \ln \mu_i + (1 - \mu_i) \ln(1 - \mu_i) \right]
  \label{eq:huang}
\end{equation}

donde $\mu_i$ es la función de membresía difusa del nivel de gris $i$ respecto a la
clase más cercana. La intuición es que el umbral óptimo es aquel que hace más \emph{clara}
la pertenencia de cada vóxel a su clase, reduciendo al mínimo la ambigüedad global.

A diferencia de Otsu, Huang no asume que las clases tengan distribuciones gaussianas ni
frecuencias similares. Su función de costo difuso le confiere mayor robustez ante
histogramas con modos de distinta amplitud o clases no gaussianas. Tiende a producir
umbrales más conservadores (más bajos) que Otsu, capturando mayor detalle en la clase
de menor intensidad.

\subsection{Método de Triangle}

El método de Triangle \cite{zack1977} opera sobre la forma geométrica del histograma.
El algoritmo traza una línea recta entre el pico máximo del histograma y su extremo más
alejado, y busca el punto del histograma que maximiza la distancia perpendicular a dicha
línea:

\begin{equation}
  T^* = \arg\max_{T}\; d_\perp\!\left(H(T),\, \overline{P_{\max}P_{\text{fin}}}\right)
  \label{eq:triangle}
\end{equation}

donde $H(T)$ es el valor del histograma en $T$, y $\overline{P_{\max}P_{\text{fin}}}$ es
la línea que une el pico máximo con el extremo del histograma. El método asume que el
histograma tiene un \textbf{pico dominante} con una cola de baja frecuencia que representa
la clase minoritaria. Es especialmente efectivo cuando una clase es muy pequeña respecto
a la otra (histograma fuertemente asimétrico), situación común en imágenes donde el fondo
domina el campo de visión.

\begin{table}[H]
\centering
\renewcommand{\arraystretch}{1.5}
\begin{tabularx}{\textwidth}{lXXX}
\toprule
\textbf{Aspecto} & \textbf{Otsu} & \textbf{Huang} & \textbf{Triangle} \\
\midrule
Criterio &
  Máxima varianza entre clases &
  Mínima entropía difusa &
  Máxima distancia geométrica \\
Supuesto del histograma &
  Bimodal, clases similares &
  Bimodal, más flexible &
  Pico dominante con cola \\
Distribución de clases &
  Gaussianas, frecuencias similares &
  No gaussianas toleradas &
  Clase dominante + clase minoritaria \\
Tendencia del umbral &
  Valle del histograma &
  Umbral conservador (bajo) &
  Umbral alto (excluye la cola) \\
Fortaleza &
  Histogramas balanceados &
  Histogramas asimétricos o solapados &
  Histogramas fuertemente asimétricos \\
Debilidad &
  Clases muy desiguales en tamaño &
  Sobreestima la clase minoritaria &
  Falla en histogramas bimodales balanceados \\
\bottomrule
\end{tabularx}
\caption{Comparación conceptual de los tres métodos de umbralización automática.}
\label{tab:metodos}
\end{table}

% ============================================================
\section{Implementación}
% ============================================================

Los tres filtros de umbralización se implementaron en Python utilizando la biblioteca ITK,
siguiendo el mismo pipeline de entrada/salida:

\begin{tcolorbox}[infostyle, title={Pipeline de segmentación por umbralización (los tres métodos)}]
\textbf{Entrada:} Imagen médica en formato NIfTI (\texttt{.nii} / \texttt{.nii.gz})\\[0.3em]
\textbf{Paso 1 — Lectura ITK:} \texttt{itk.imread()} carga el volumen 3D.\\[0.3em]
\textbf{Paso 2 — Umbralización:} se instancia el filtro correspondiente
(\texttt{OtsuThresholdImageFilter}, \texttt{HuangThresholdImageFilter} o
\texttt{TriangleThresholdImageFilter}) con \texttt{InsideValue=0} y
\texttt{OutsideValue=255}.\\[0.3em]
\textbf{Paso 3 — Recuperación del umbral:} \texttt{GetThreshold()} retorna el valor $T^*$
calculado automáticamente.\\[0.3em]
\textbf{Paso 4 — Extracción del corte central:} se extrae el corte axial central del
volumen segmentado para visualización.\\[0.3em]
\textbf{Salida:} Máscara binaria NIfTI + PNG de comparación (Original | Otsu | Huang | Triangle).
\end{tcolorbox}

Los tres filtros comparten la misma interfaz en ITK: reciben una imagen de entrada, calculan
el umbral automáticamente a partir del histograma interno (\texttt{NumberOfHistogramBins=128}
por defecto) y producen una imagen de salida binaria donde los vóxeles con intensidad inferior
al umbral toman el valor \texttt{InsideValue} y los superiores el valor \texttt{OutsideValue}.
La convención usada en este taller asigna \texttt{InsideValue=0} (negro, fondo) y
\texttt{OutsideValue=255} (blanco, objeto), de modo que las regiones de mayor intensidad
en la imagen original aparecen en blanco en la máscara.

% ============================================================
\section{Dataset — Imágenes utilizadas}
% ============================================================

Los experimentos utilizan imágenes representativas de tres modalidades de imagen médica,
seleccionadas para abarcar un amplio rango de características de histograma.

\begin{table}[H]
\centering
\renewcommand{\arraystretch}{1.3}
\begin{tabularx}{\textwidth}{llXX}
\toprule
\textbf{Imagen} & \textbf{Modalidad} & \textbf{Estructura de interés} & \textbf{Característica del histograma} \\
\midrule
\texttt{coronacases\_org\_001} &
  CT de tórax &
  Pulmones, tejidos blandos, estructuras óseas &
  Multimodal: pico estrecho de aire (HU bajos) y distribución amplia de tejidos \\
\texttt{BraTS-GLI-00000-000-t1c} &
  MRI T1c cerebral &
  Tumor cerebral (glioma), materia blanca, materia gris &
  Bimodal: fondo negro (máscara craneal) y tejido cerebral con realce tumoral \\
\texttt{04\_01} &
  Ultrasonido endorrectal (US) &
  Próstata, tejido periprostático &
  Histograma amplio y distribuido por ruido speckle \\
\bottomrule
\end{tabularx}
\caption{Imágenes médicas utilizadas en los experimentos de umbralización.}
\end{table}

% ============================================================
\section{Resultados}
% ============================================================

\subsection{CT de tórax — \texttt{coronacases\_org\_001}}

\begin{figure}[H]
  \centering
  \includegraphics[width=\textwidth]{output/comparisons/coronacases_org_001_comparison.png}
  \caption{CT de tórax: Original, Otsu, Huang y Triangle (corte axial central).
           Las regiones blancas corresponden a los vóxeles con intensidad superior al
           umbral calculado por cada método.}
  \label{fig:ct_thorax}
\end{figure}

El CT de tórax presenta el histograma más favorable para la umbralización global: existe
una separación clara entre el aire intrapulmonar (intensidades bajas en Hounsfield) y
los tejidos blandos y óseos (intensidades altas). Los tres métodos producen máscaras
binarias reconocibles, aunque con diferencias notables en el umbral elegido.

\textbf{Otsu} ubica el umbral en el valle del histograma entre los dos modos principales,
generando una segmentación limpia donde los pulmones (clase de baja intensidad) quedan
en negro y el resto del tórax —tejidos blandos, costillas y tejido subcutáneo— en blanco.
El resultado es una máscara de alta legibilidad diagnóstica con bordes bien definidos.

\textbf{Huang}, al minimizar la entropía difusa, elige un umbral
más bajo que Otsu, capturando también estructuras pulmonares internas de intensidad
intermedia: vasos, bronquios y nódulos. Las regiones blancas son más fragmentadas,
lo que puede ser ventajoso para análisis detallado del parénquima pulmonar pero
introduce más falsos positivos en la máscara.

\textbf{Triangle} produce un umbral intermedio, aunque más cercano a Otsu que a Huang
en esta imagen. La máscara resultante conserva las grandes estructuras torácicas con
menor fragmentación interna que Huang, lo que refleja la sensibilidad del método de
Triangle a la forma asimétrica del histograma CT.

\subsection{MRI cerebral T1c — \texttt{BraTS-GLI-00000-000-t1c}}

\begin{figure}[H]
  \centering
  \includegraphics[width=\textwidth]{output/comparisons/BraTS-GLI-00000-000-t1c_comparison.png}
  \caption{MRI T1c cerebral (BraTS): Original, Otsu, Huang y Triangle (corte axial central).
           La imagen original ya presenta máscara craneal (fondo negro). Las regiones
           blancas en las máscaras corresponden a tejido cerebral de alta intensidad y
           realce tumoral por contraste con gadolinio.}
  \label{fig:mri_brain}
\end{figure}

La imagen MRI T1c es una secuencia con contraste de gadolinio, donde el tejido tumoral
activo (glioma de alto grado) se realza y aparece con mayor intensidad que el tejido
cerebral sano circundante. La imagen ya cuenta con una máscara de remoción de cráneo
aplicada, por lo que el fondo es negro puro, lo que simplifica la estructura del
histograma.

\textbf{Otsu} genera la segmentación con mayor detalle estructural de los tres métodos,
preservando la materia blanca, la materia gris y las zonas de realce tumoral en la
máscara blanca. El umbral resulta relativamente bajo porque Otsu balancea las dos clases
(fondo negro y tejido cerebral) y el tejido ocupa una fracción significativa del volumen.

\textbf{Huang} produce un resultado visualmente similar a Otsu pero con el umbral
levemente desplazado hacia valores superiores de intensidad. Esto resulta en una máscara
ligeramente más fragmentada, donde algunas regiones de tejido de intensidad intermedia
(materia gris más oscura) son excluidas. Las zonas de realce tumoral, al ser las más
brillantes de la imagen, permanecen en blanco en ambos métodos.

\textbf{Triangle} es el método más conservador en esta modalidad: su umbral es el más
alto de los tres, por lo que solo los tejidos de mayor intensidad —materia blanca muy
brillante y realce tumoral— quedan en la máscara. El resultado es una máscara más
restringida que podría ser de utilidad para aislar específicamente las regiones de
realce por contraste, pero a costa de excluir tejido cerebral normal de intensidad
moderada.

\begin{table}[H]
\centering
\renewcommand{\arraystretch}{1.4}
\begin{tabularx}{\textwidth}{lXX}
\toprule
\textbf{Método} & \textbf{Tejido capturado (blanco)} & \textbf{Observación clínica} \\
\midrule
Otsu &
  Materia blanca + materia gris + realce tumoral &
  Segmentación más completa del parénquima cerebral \\
Huang &
  Materia blanca + realce tumoral + parte de la materia gris &
  Umbral intermedio; mayor fragmentación en zonas de baja señal \\
Triangle &
  Materia blanca brillante + realce tumoral principalmente &
  Útil para aislar realce por contraste; excluye materia gris \\
\bottomrule
\end{tabularx}
\caption{Comparación de la región segmentada por cada método en la MRI T1c.}
\end{table}

\subsection{Ultrasonido endorrectal — \texttt{04\_01}}

\begin{figure}[H]
  \centering
  \includegraphics[width=\textwidth]{output/comparisons/04_01_comparison.png}
  \caption{Ultrasonido endorrectal de próstata: Original, Otsu, Huang y Triangle
           (corte axial central). El ruido speckle característico del ultrasonido
           distribuye las intensidades de forma continua, dificultando la umbralización.}
  \label{fig:us_prostate}
\end{figure}

El ultrasonido es la modalidad más desafiante para la umbralización global. El ruido
\textit{speckle} inherente a la adquisición ultrasónica genera intensidades distribuidas
de forma continua en toda la imagen, sin producir los modos bien diferenciados que
sustentan la asunción bimodal de estos métodos. Como consecuencia, los tres algoritmos
producen máscaras fragmentadas con alta densidad de falsos positivos y falsos negativos.

\textbf{Otsu} genera una segmentación muy fragmentada donde las regiones blancas y negras
se distribuyen de forma caótica en toda la región prostática. El histograma del ultrasonido
es aproximadamente unimodal con cola extendida; la maximización de la varianza entre clases
ubica el umbral en una zona del histograma que no corresponde a una frontera anatómica real.

\textbf{Huang} produce un resultado muy similar a Otsu en términos de fragmentación.
La entropía difusa no encuentra una separación clara en el histograma continuo del
ultrasonido, lo que resulta en una máscara igualmente ruidosa. La similitud entre ambos
métodos refleja que cuando el histograma no es bimodal, tanto el criterio de varianza
(Otsu) como el de entropía (Huang) convergen a umbrales de calidad comparable pero
insatisfactoria.

\textbf{Triangle} es el que produce la segmentación visualmente más diferenciada de
los tres para esta modalidad, concentrando las regiones blancas en una franja central de
la imagen que se aproxima parcialmente a la región prostática. Esto se debe a que el
histograma del ultrasonido tiene un pico de baja intensidad (fondo oscuro) con una cola
hacia valores medios, estructura que se ajusta mejor al supuesto geométrico del método
de Triangle. Sin embargo, el resultado sigue siendo inaceptable para uso diagnóstico,
lo que confirma que la umbralización global no es el método de elección para imágenes
de ultrasonido.

% ============================================================
\section{Efecto del número de barras del histograma (\texttt{NumberOfHistogramBins})}
% ============================================================

El parámetro \texttt{NumberOfHistogramBins} ($b$) controla la resolución con la que los
filtros de umbralización automática de ITK discretizan el rango de intensidades de la
imagen antes de aplicar su criterio de optimización. Un valor bajo de $b$ \textbf{suaviza}
el histograma agrupando intensidades distintas en el mismo bin, lo que puede simplificar
artificialmente la estructura del histograma. Un valor alto permite \textbf{mayor
resolución}, capturando variabilidad fina de intensidades pero también amplificando el
efecto del ruido sobre la estimación del umbral. El valor por defecto en ITK es
$b = 128$.

Para cuantificar este efecto se evaluaron cuatro valores: $b \in \{32, 64, 128, 256\}$,
aplicados sobre las tres modalidades de imagen. Las imágenes sin sufijo en las secciones
anteriores corresponden a $b = 128$ (valor por defecto). Los resultados se presentan a
continuación, con una subsección dedicada a cada valor de $b$.

\begin{table}[H]
\centering
\renewcommand{\arraystretch}{1.4}
\begin{tabularx}{\textwidth}{lXXX}
\toprule
\textbf{$b$} & \textbf{CT de tórax} & \textbf{MRI T1c} & \textbf{US próstata} \\
\midrule
32  & Triangle converge con Otsu; máscaras limpias & Sin cambio visible & Otsu/Huang más continuos; Triangle ruidoso \\
64  & Triangle comienza a separarse de Otsu & Sin cambio visible & Comportamiento intermedio \\
128 & Resultado de referencia (default ITK) & Sin cambio visible & Triangle más compacto \\
256 & Triangle degrada: tejido subcutáneo incluido; Otsu estable & Sin cambio visible & Triangle más compacto y preciso \\
\bottomrule
\end{tabularx}
\caption{Resumen del efecto de \texttt{NumberOfHistogramBins} por modalidad.}
\label{tab:bins}
\end{table}

% ------------------------------------------------------------
\subsection{$b = 32$ barras}
% ------------------------------------------------------------

Con 32 barras, el histograma queda fuertemente suavizado: cada bin agrupa un rango
amplio de intensidades, lo que tiende a \textbf{fusionar modos cercanos} y a simplificar
la estructura de la distribución. Este efecto tiene consecuencias distintas según la
modalidad y el método.

\begin{figure}[H]
  \centering
  \includegraphics[width=\textwidth]{output/comparisons/coronacases_org_001_comparison_b32.png}\\[0.8em]
  \includegraphics[width=\textwidth]{output/comparisons/BraTS-GLI-00000-000-t1c_comparison_b32.png}\\[0.8em]
  \includegraphics[width=\textwidth]{output/comparisons/04_01_comparison_b32.png}
  \caption{Comparación con $b = 32$ barras. De arriba a abajo: CT de tórax, MRI T1c
           cerebral, ultrasonido endorrectal. Columnas: Original, Otsu, Huang, Triangle.}
  \label{fig:b32}
\end{figure}

\textbf{CT de tórax:} Con $b = 32$, el método Triangle produce su resultado más similar
a Otsu de toda la serie experimental. El histograma suavizado reduce los picos secundarios
de intensidades intermedias (tejido subcutáneo, estructuras mediastínicas), lo que lleva
a Triangle a identificar un umbral más conservador y cercano al valle principal del
histograma. El resultado es una máscara limpia donde los pulmones quedan bien delimitados
sin el ruido periférico que aparece con resoluciones más altas. \textbf{Otsu} y
\textbf{Huang} son prácticamente insensibles a este cambio: sus criterios globales (varianza
entre clases y entropía difusa) se estabilizan incluso con resolución baja de histograma.

\textbf{MRI T1c cerebral:} Los tres métodos producen resultados visualmente idénticos a
los obtenidos con $b = 128$. La estructura bimodal de esta imagen (fondo negro puro +
tejido cerebral) es tan pronunciada que incluso con solo 32 bins todos los algoritmos
identifican el mismo valle del histograma y calculan umbrales equivalentes. Esta
insensibilidad es una señal positiva: indica que el histograma de la imagen tiene una
separación de modos robusta y que la segmentación es estable independientemente del valor
de $b$ elegido.

\textbf{Ultrasonido:} Con $b = 32$ los métodos Otsu y Huang producen máscaras con regiones
blancas ligeramente más amplias y continuas que con resoluciones superiores. El histograma
suavizado atenúa la variabilidad local del speckle, lo que empuja el umbral hacia un valor
que captura más píxeles como objeto. Triangle en cambio presenta mayor fragmentación a
$b = 32$ que a valores superiores, ya que la cola del histograma queda mal definida con
tan pocos bins y el algoritmo geométrico pierde precisión. En ningún caso el resultado
es clínicamente aceptable, pero la diferencia entre métodos es observable.

\textbf{Recomendación de uso:} $b = 32$ puede ser una opción razonable cuando se trabaja
con imágenes CT y se desea que Triangle se comporte de forma más conservadora y parecida
a Otsu, o cuando se busca estabilidad del umbral ante ruido leve en el histograma.
No es recomendable para imágenes de ultrasonido ni para aplicaciones que requieran
precisión fina en la detección del umbral.

% ------------------------------------------------------------
\subsection{$b = 64$ barras}
% ------------------------------------------------------------

\begin{figure}[H]
  \centering
  \includegraphics[width=\textwidth]{output/comparisons/coronacases_org_001_comparison_b64.png}\\[0.8em]
  \includegraphics[width=\textwidth]{output/comparisons/BraTS-GLI-00000-000-t1c_comparison_b64.png}\\[0.8em]
  \includegraphics[width=\textwidth]{output/comparisons/04_01_comparison_b64.png}
  \caption{Comparación con $b = 64$ barras. De arriba a abajo: CT de tórax, MRI T1c
           cerebral, ultrasonido endorrectal.}
  \label{fig:b64}
\end{figure}

% [Subsección a cargo del compañero correspondiente]
\textit{Análisis pendiente de completar por integrante del equipo.}

% ------------------------------------------------------------
\subsection{$b = 128$ barras (valor de referencia)}
% ------------------------------------------------------------

\begin{figure}[H]
  \centering
  \includegraphics[width=\textwidth]{output/comparisons/coronacases_org_001_comparison_b128.png}\\[0.8em]
  \includegraphics[width=\textwidth]{output/comparisons/BraTS-GLI-00000-000-t1c_comparison_b128.png}\\[0.8em]
  \includegraphics[width=\textwidth]{output/comparisons/04_01_comparison_b128.png}
  \caption{Comparación con $b = 128$ barras (valor por defecto de ITK). De arriba a abajo:
           CT de tórax, MRI T1c cerebral, ultrasonido endorrectal.}
  \label{fig:b128}
\end{figure}

% [Subsección a cargo del compañero correspondiente]
\textit{Análisis pendiente de completar por integrante del equipo.}
% ------------------------------------------------------------
\subsection{$b = 256$ barras}
% ------------------------------------------------------------
Con 256 barras, el histograma queda bien poblado, en tanto tiene asignada cada intensidad 1-1, esto es, una barra (bin) del histograma es una nota de intensidad. 
esto tiende a \textbf{resaltar el ruido} y a amplificar
la estructura de la distribución con su ruido corresponiente. Este efecto tiene consecuencias distintas según la
modalidad y el método.

\begin{figure}[H]
  \centering
  \includegraphics[width=\textwidth]{output/comparisons/coronacases_org_001_comparison_b256.png}\\[0.8em]
  \includegraphics[width=\textwidth]{output/comparisons/BraTS-GLI-00000-000-t1c_comparison_b256.png}\\[0.8em]
  \includegraphics[width=\textwidth]{output/comparisons/04_01_comparison_b256.png}
  \caption{Comparación con $b = 256$ barras. De arriba a abajo: CT de tórax, MRI T1c
           cerebral, ultrasonido endorrectal.}
  \label{fig:b256}
\end{figure}

\textbf{CT de tórax:} Con $b = 256$, el método Triangle produce un resultado muy poco útil, en tanto el cálculo del umbral
en éste método depende de la cantidad y altura de bins en el histograma, lo que lo hace susceptible en éste caso de no producir buenos umbrales.
Se notan algunas siluetas de los bordes que permiten mantener una idea ligera de la imagen pero se piederde más información de la deseada. 
\textbf{Otsu} y
\textbf{Huang} mantienen una estructura similar, aunque en éste caso, Otsu tiende a uniformizar mejor los blancos de la región de interés mientras que
Huang logra mantener un contraste a pesar de que opaca algunas regiones de forma muy fuerte.

\textbf{MRI T1c cerebral:} Los tres métodos producen resultados visualmente idénticos a
los obtenidos con $b = 128$. La estructura bimodal de esta imagen (fondo negro puro +
tejido cerebral) es tan pronunciada que, dados los 256 bins, no se logra tener un detalle gráfico de la imagen que sea diferenciador.
En éste caso, el uso de 256 está muy poco recomendado.

\textbf{Ultrasonido:} Con $b = 256$ los métodos Otsu y Huang producen máscaras con regiones
marcadas, pero con sutilezas: Otsu resalta parte del ruido mientras qeu Huang opaca recgiones de posible interés. 
El histograma detallado enfativza la variabilidad local, lo que explica la variabilidad en los resultados.
Triangle en cambio, nuevamente, dada la sensiblidad al tamaño del histograma, no suele ser tan útil.

\textbf{Recomendación de uso:} $b = 256$ puede ser una opción razonable cuando hay un requerimiento
de exactitud para MRI o para CT. Por otro lado, su uso en ecografías está poco recomendado pues puede 
ya sea acabar con regiones de interés o aumentar el ruido.

% ============================================================
\section{Análisis Comparativo}
% ============================================================

\subsection{Comparación entre los tres métodos de umbralización}

\begin{table}[H]
\centering
\renewcommand{\arraystretch}{1.5}
\begin{tabularx}{\textwidth}{lXXX}
\toprule
\textbf{Aspecto} &
  \textbf{Otsu} &
  \textbf{Huang} &
  \textbf{Triangle} \\
\midrule
Umbral en CT tórax &
  Valle del histograma, limpio &
  Más bajo; captura estructuras internas &
  Intermedio; similar a Otsu \\
Umbral en MRI T1c &
  Bajo; captura todo el parénquima &
  Intermedio; más fragmentado &
  Alto; solo tejido muy brillante \\
Umbral en US próstata &
  Insatisfactorio (speckle) &
  Insatisfactorio (similar a Otsu) &
  Parcialmente mejor; concentra región central \\
Sensibilidad al histograma &
  Alta: requiere bimodalidad clara &
  Moderada: tolera asimetría &
  Alta: asume pico dominante \\
Adecuación clínica CT &
  Excelente &
  Buena (mayor detalle) &
  Buena \\
Adecuación clínica MRI &
  Muy buena (segmentación completa) &
  Buena (realce tumoral) &
  Moderada (solo alta intensidad) \\
Adecuación clínica US &
  Pobre &
  Pobre &
  Parcial (inaceptable para diagnóstico) \\
Sensibilidad a \texttt{NumberOfHistogramBins} &
  Baja: estable en todos los valores de $b$ &
  Baja: estable en todos los valores de $b$ &
  Alta en CT: $b=32$ converge con Otsu; $b=256$ degrada \\
\bottomrule
\end{tabularx}
\caption{Comparación del desempeño de los tres métodos de umbralización por modalidad.}
\label{tab:comparacion}
\end{table}

\subsection{Relación entre estructura del histograma y calidad de segmentación}

Los resultados experimentales confirman el principio conceptual expuesto en clase
\cite{bankman2009}: la calidad de la segmentación por umbralización global está
directamente condicionada por la estructura bimodal del histograma. En la CT de tórax,
donde existe una separación clara entre el aire y los tejidos blandos, los tres métodos
producen segmentaciones reconocibles y clínicamente interpretables. En la MRI cerebral
T1c, la presencia de una máscara de fondo negro produce un histograma con dos regiones
diferenciadas, lo que favorece la umbralización aunque con diferencias en el nivel de
detalle capturado. En el ultrasonido, donde el ruido speckle destruye la bimodalidad
del histograma, ningún método logra una segmentación satisfactoria.

Esta dependencia de la distribución de intensidades es la limitación fundamental de la
umbralización global como método de segmentación, y motiva el uso de métodos más
sofisticados —crecimiento de regiones, modelos de contornos activos, campos aleatorios
de Markov— para modalidades con histogramas complejos o en presencia de ruido.

\subsection{Implicaciones para la segmentación de tumores cerebrales (proyecto)}

En el contexto del proyecto de segmentación de gliomas con el dataset BraTS, la elección
del método de umbralización influye directamente en la calidad de la máscara inicial.
Otsu captura el tejido cerebral completo y puede usarse como paso previo para aislar
el cerebro del fondo. Triangle, al producir un umbral más alto, es útil para obtener
una máscara restringida al realce tumoral en secuencias T1c, que puede servir como
región semilla para métodos de crecimiento de regiones. Huang, con su comportamiento
intermedio, ofrece un balance entre sensibilidad y especificidad en tejidos con
intensidades solapadas.

% ============================================================
\section{Conclusiones}
% ============================================================

\begin{enumerate}

  \item \textbf{Método de Otsu:} Produce los mejores resultados generales en imágenes
    con histograma bimodal balanceado. En CT de tórax ofrece una separación limpia entre
    pulmones y tejidos blandos; en MRI cerebral captura el parénquima completo. Su
    desempeño es consistente con la teoría: al maximizar la varianza entre clases,
    opera óptimamente cuando los dos grupos tienen frecuencias de aparición similares.
    En ultrasonido falla por la ausencia de bimodalidad real en el histograma.

  \item \textbf{Método de Huang:} El criterio de entropía difusa le confiere mayor
    robustez ante histogramas con modos asimétricos o intensidades solapadas, produciendo
    umbrales más conservadores que capturan mayor detalle en la clase de menor intensidad.
    En CT de tórax esto se traduce en mayor captura de estructuras pulmonares internas;
    en MRI cerebral en mayor fragmentación. Para imágenes de ultrasonido su desempeño
    es equivalente al de Otsu, confirmando que el problema reside en la distribución
    de intensidades y no en el criterio de optimización.

  \item \textbf{Método de Triangle:} Su supuesto de un pico dominante con una cola
    minoritaria lo hace especialmente adecuado para imágenes donde la clase de interés
    es pequeña respecto al fondo. En MRI T1c es útil para aislar el realce tumoral,
    la región de mayor intensidad. En ultrasonido logra parcialmente separar el centro
    de la imagen, aunque el resultado sigue siendo insatisfactorio para uso clínico.

  \item \textbf{Limitación fundamental:} Los tres métodos comparten la misma restricción
    estructural: asumen un histograma bimodal. Cuando esta condición no se cumple
    —como en imágenes de ultrasonido con ruido speckle— ningún método de umbralización
    global produce una segmentación clínicamente útil. En estos casos es necesario
    recurrir a métodos basados en regiones, contornos o modelos probabilísticos.

  \item \textbf{Rol de ITK:} La biblioteca ITK proporciona implementaciones eficientes
    y correctas de los tres métodos a través de \texttt{OtsuThresholdImageFilter},
    \texttt{HuangThresholdImageFilter} y \texttt{TriangleThresholdImageFilter}, con
    interfaz uniforme que facilita la comparación sistemática. El cálculo automático
    del histograma interno y la recuperación del umbral con \texttt{GetThreshold()}
    permiten una implementación concisa y reproducible sobre volúmenes 3D.

  \item \textbf{Efecto de \texttt{NumberOfHistogramBins}:} La resolución del histograma
    interno afecta de forma desigual a los tres métodos. Otsu y Huang son robustos ante
    variaciones de $b$: sus criterios globales convergen al mismo umbral independientemente
    de si se usan 32 o 256 bins. Triangle es el más sensible: en CT de tórax, $b = 32$
    produce un umbral más conservador y similar a Otsu (deseable), mientras que $b = 256$
    incorpora estructuras periféricas no deseadas. En MRI T1c con máscara craneal ningún
    valor de $b$ altera el resultado, lo que confirma la solidez de la bimodalidad de esa
    imagen. En ultrasonido el efecto es moderado y no supera la limitación estructural del
    speckle. El valor por defecto $b = 128$ es un balance adecuado para la mayoría de
    escenarios clínicos.

\end{enumerate}

% ============================================================
\section*{Referencias}
\addcontentsline{toc}{section}{Referencias}
% ============================================================

\begin{thebibliography}{9}

\bibitem{bankman2009}
  Bankman, I.\,N. (Ed.) (2009).
  \textit{Handbook of Medical Image Processing and Analysis} (2nd ed.).
  Academic Press.

\bibitem{otsu1979}
  Otsu, N. (1979).
  A threshold selection method from gray-level histograms.
  \textit{IEEE Transactions on Systems, Man, and Cybernetics}, 9(1), 62--66.
  \url{https://doi.org/10.1109/TSMC.1979.4310076}

\bibitem{huang1995}
  Huang, L.-K., \& Wang, M.-J. (1995).
  Image thresholding by minimizing the measures of fuzziness.
  \textit{Pattern Recognition}, 28(1), 41--51.
  \url{https://doi.org/10.1016/0031-3203(94)E0043-K}

\bibitem{zack1977}
  Zack, G.\,W., Rogers, W.\,E., \& Latt, S.\,A. (1977).
  Automatic measurement of sister chromatid exchange frequency.
  \textit{Journal of Histochemistry and Cytochemistry}, 25(7), 741--753.
  \url{https://doi.org/10.1177/25.7.70454}

\bibitem{itk}
  ITK Development Team. (2024).
  \textit{Insight Toolkit (ITK) v5.4}.
  \url{https://itk.org}

\end{thebibliography}

\end{document}
